\documentclass[11pt,a4paper]{article}
\usepackage[T1]{fontenc}
\usepackage{lmodern}
\usepackage[margin=26mm]{geometry}
\usepackage{amsmath,amssymb,amsthm}
\usepackage{microtype}
\usepackage{enumitem}
\usepackage[hidelinks]{hyperref}
\setlength{\parskip}{3pt}
\setlength{\emergencystretch}{3em}
\newtheorem{theorem}{Theorem}[section]
\newtheorem{lemma}[theorem]{Lemma}
\newtheorem{corollary}[theorem]{Corollary}
\theoremstyle{remark}
\newtheorem{remark}[theorem]{Remark}
\newcommand{\Z}{\mathbb Z}
\newcommand{\N}{\mathbb N}
\newcommand{\Q}{\mathbb Q}
\newcommand{\F}{\mathbb F}
\newcommand{\ps}[1]{#1[[x]]}
\hypersetup{pdftitle={A306082: Square-partition Stirling periodicity},pdfauthor={Dev Goyal}}
\title{A306082\\[3pt]Square-partition Stirling periodicity}
\author{Dev Goyal}
\date{24 October 2026}
\begin{document}
\maketitle
\begin{center}
\fbox{\begin{minipage}{0.91\linewidth}\small
\textbf{AI contribution and review status.}
ChatGPT (OpenAI) generated the mathematical exposition, proof development,
Lean interfaces, and checks for this collection in response to Dev Goyal's
prompts. Dev Goyal is the sole human author. Independent human review is
pending. Original sequence authors, known mathematical results, and adapted
upstream code retain their separate credits. This document makes no claim
of publication, OEIS acceptance, or worldwide priority.
\end{minipage}}
\end{center}
\begin{abstract}
For the explicitly defined factorial Stirling transform corresponding to A306082, we give a self-contained proof of eventual periodicity modulo every positive integer. The tail beginning at the modulus has period dividing Euler's totient. This is a formalization and specialization of known mathematics, credited to the prior work below.
\end{abstract}
\section{Exact sequence and result}
Let $S(n,k)$ be the Stirling numbers of the second kind. Their initial
values are $S(0,0)=1$, $S(0,k)=0$ for $k>0$, and $S(n,0)=0$ for
$n>0$; their recurrence is
\[
S(n+1,k+1)=(k+1)S(n,k+1)+S(n,k).
\]
The finite source formula for A306082 \cite{oeis} is
\[
a_n=\sum_{k=0}^n c_k k!S(n,k),\qquad c_k=p_{\mathrm{sq}}(k).
\]
The source offset is zero. Here $p_{\mathrm{sq}}(k)$ counts partitions of $k$ whose positive parts are all perfect squares; the empty partition has count one. The integral generating series is $G(t)=\prod_{j\ge1}(1-t^{j^2})^{-1}$.

For every $m\ge1$ and $n\ge m$, we prove
$a_{n+\varphi(m)}\equiv a_n\pmod m$. We also prove that the least
positive period of $t\mapsto a_{m+t}\bmod m$ exists, is itself a
period, and divides $\varphi(m)$.
\section{The universal finite-transform theorem}
For an arbitrary sequence of integers $(c_k)_{k\ge0}$, write
\[
T_c(n)=\sum_{k=0}^n c_k k!S(n,k).
\]
No growth, sign, positivity, or recurrence hypothesis is needed.
\begin{theorem}\label{thm:period}
For every $m\ge1$ and $n\ge m$,
\[
T_c(n+\varphi(m))\equiv T_c(n)\pmod m.
\]
The tail $t\mapsto T_c(m+t)\bmod m$ has a least positive period $d$;
this $d$ is itself a period and divides $\varphi(m)$.
\end{theorem}
\begin{proof}
If $m=1$ every residue is zero, so the assertions hold with $d=1$.
Suppose $m>1$. First, for every natural number $a$ and $n\ge m$,
\begin{equation}\label{eq:powers}
a^{n+\varphi(m)}\equiv a^n\pmod m.
\end{equation}
Check each maximal prime-power factor $p^e$ of $m$. If $p\mid a$,
then $e\le p^e\le m\le n$, so $p^e$ divides both powers. If
$p\nmid a$, Euler's theorem gives period $\varphi(p^e)$ for powers
of $a$ modulo $p^e$. The prime-power product formula for $\varphi$
shows that $\varphi(p^e)\mid\varphi(m)$. Thus the congruence holds
modulo every prime-power factor, and hence modulo $m$.

The Stirling recurrence gives the falling-factorial expansion
\begin{equation}\label{eq:triangular-power}
a^n=\sum_{k=0}^{a}\binom ak k!S(n,k)\qquad(a\le n).
\end{equation}
For completeness, it follows by counting all functions from an
$n$-element set to an $a$-element set: their $k$ nonempty fibers
partition the domain in $S(n,k)$ ways, and the fibers may be assigned
distinct target values in $a(a-1)\cdots(a-k+1)=\binom ak k!$ ways.
Equivalently, both sides satisfy the induction obtained from the
Stirling recurrence, including $n=0$.

Fix $n\ge m$. Strong induction on $k<m$ now proves
\[
k!S(n+\varphi(m),k)\equiv k!S(n,k)\pmod m.
\]
Apply \eqref{eq:triangular-power} with $a=k$ to the two exponents.
Their left sides are congruent by \eqref{eq:powers}. Each term with
index below $k$ cancels by induction; the remaining multiplier is
$\binom kk=1$. No division by $k!$ modulo $m$ is involved.

For every $k\ge m$, the factor $k!$ is divisible by $m$. The two
transforms therefore reduce to the same fixed finite range $k<m$;
the columnwise congruence just proved gives their equality modulo $m$.

Since $\varphi(m)>0$, at least one positive tail period exists.
Let $d$ be the least such period and let $p$ be any other period.
Write $p=qd+r$ with $0\le r<d$. If $f(t)=T_c(m+t)\bmod m$, then
$f(t+r)=f(t+r+qd)=f(t+p)=f(t)$ for all $t\ge0$.
Thus $r$ is a period, and minimality forces $r=0$.
Taking $p=\varphi(m)$ proves $d\mid\varphi(m)$.
\end{proof}

\section{Generating functions and the application}
The preceding theorem directly applies to the finite formula defining
the sequence in this document. It proves both its totient period and
the existence and divisibility property of its least tail period.

The usual exponential-generating-function interpretation is also
mathematically valid. In $\Q[[x]]$ one has
\[
(e^x-1)^k=k!\sum_{n\ge k}S(n,k)\frac{x^n}{n!}.
\]
Indeed, expanding $(e^x-1)^k$ coefficientwise counts ways to distribute
$n$ labeled elements among $k$ labeled nonempty boxes. These are
surjections, whose number is $k!S(n,k)$. At any fixed degree all
coefficient operations are finite. Consequently, for
$G(t)=\sum_{k\ge0}c_kt^k$,
\[
G(e^x-1)=\sum_{n\ge0}T_c(n)\frac{x^n}{n!}.
\]
This explanation identifies the source's generating-function description
with its finite transform. \textbf{The formal-exponential substitution
identity itself is not separately formalized in this collection.}
The Lean theorems instead start from the explicit finite Stirling sum,
with the exact integer weight specified above. No unproved EGF bridge is
used as a Lean axiom or hypothesis.

\section{Prior result and attribution}
The periodicity is known mathematics, not a newly solved open problem.
Danesh's Theorem 4.1 and Corollary 4.2 \cite{danesh} already give a
stronger universal eventual Carmichael period $\lambda(m)$ for integral
substitutions $G(e^x-1)$. This covers the present totient-period
conclusion as an eventual-period statement. The explicit threshold
$n\ge m$ used here is sufficient; optimality or priority for that
threshold is not asserted. Earlier Fubini and $r$-Fubini work by
Schreyer is also relevant \cite{schreyer}. The formal result here
proves the totient bound; it does not assert a Lean proof of the
stronger Carmichael bound.

The Stirling recurrence and part of the power-expansion implementation
were adapted from Apache-2.0 Mathlib source \cite{mathlib}.
The original copyright, authors, license, and change notice are retained.
The modular transfer and named specialization proofs were generated
for this project. This AI disclosure does not claim authorship of
the adapted upstream code.
\section{Lean proof and exact scope}
The dedicated entry file is \path{lean/Entries/A306082.lean}. Its declarations lie in namespace \texttt{FinalCollection.A306082}. The principal checked endpoints are:
\begin{itemize}[nosep]
\item \texttt{source\_formula}: the actual finite sum with the stated weight.
\item \texttt{period}: the unbounded congruence for every $m\ge1$ and $n\ge m$.
\item \texttt{least\_period\_divides}: divisibility of the least positive tail period.
\item \texttt{least\_period\_specification}: positivity, periodicity, divisibility, and minimality together.
\item \texttt{square\_weight\_definition}: the finite cardinality of partitions with square parts.
\end{itemize}
The full implementations are in \path{lean/More/StirlingPeriodicity.lean}.
This entry formalizes a known periodicity result with an explicit sufficient onset. The actual finite integer weight and sum are included in the source formula; no undefined external sequence is assumed to agree with the OEIS entry. The EGF identity is proved in prose above but is not a separate Lean endpoint. 

The companion is a separate, explicitly typed Lean entry file. Its proofs
invoke the preserved implementation modules; the complete dependency source
is included in the same archive. Shared lemmas are reused rather than
assumed. The entry file is not a self-contained replacement for Mathlib or
the supplied common library.

The project pins Lean 4.19.0 and Mathlib revision
\begin{center}\small\ttfamily
c44e0c8ee63ca166450922a373c7409c5d26b00b.
\end{center}
The collection checker recompiles the proof sources and audits the
theorems' transitive axiom dependencies. Only \texttt{propext},
\texttt{Classical.choice}, and \texttt{Quot.sound} are permitted.
The delivered verification reports record the actual successful runs.
Finite numerical checks are corroboration; the displayed conclusions are
unbounded theorems. An AI review of source correspondence does not replace
independent human review.
\begin{thebibliography}{9}
\raggedright
\bibitem{oeis} OEIS Foundation, entry A306082,
\url{https://oeis.org/A306082}. Source checked 24 October 2026.
\bibitem{danesh} P. Danesh, \emph{Eventual periodicity of weighted
Stirling transforms and two conjectures of Bala}, version 1,
26 August 2026, Theorem 4.1 and Corollary 4.2.
\url{https://doi.org/10.33774/coe-2026-4d6jr}.
\bibitem{schreyer} B. Schreyer, \emph{Rigged Horse Numbers and their
Modular Periodicity}, arXiv:2409.03799v2, 3 December 2024.
\url{https://arxiv.org/abs/2409.03799}.
\bibitem{mathlib} Mathlib contributors. The adapted Stirling source is
\texttt{Mathlib.Combinatorics.Enumerative.Stirling} at revision
\path{4ed598f443601d8b5abccc2a14e3ba408ae47056}, copyright 2025 Beibei Xiong;
authors Beibei Xiong, Yu Shao, Weijie Jiang, and Zhengfeng Yang.
Apache-2.0 attribution and license are preserved in
\path{third_party/mathlib/NOTICE.md} and the accompanying license file.
\bibitem{repo} D. Goyal, supplied reference repository,
\url{https://github.com/dvgyl/oeis-a396843-proof}.
This separate presentation belongs to the final collected proof package;
no earlier DOI is reassigned to it.
\end{thebibliography}
\end{document}
